\chapter[Введение в RL]{Введение в обучение с подкреплением}
\label{ch:intro}

\begin{quote}
\itshape
Агент учится не потому, что ему говорят правильное действие, а открывая\,---\,через пробы и взаимодействие\,---\,какие действия приносят наибольшее долгосрочное вознаграждение.
\end{quote}

\bigskip

Обучение с подкреплением (RL) является разделом машинного обучения, посвящённым \emph{последовательному принятию решений в условиях неопределённости}.  В отличие от обучения с учителем, где заранее предоставляется фиксированный набор размеченных примеров, агент RL должен взаимодействовать со \emph{средой}, выбирать \emph{действия}, наблюдать их последствия и использовать получаемые \emph{вознаграждения} для улучшения своего поведения со временем.

Данная глава знакомит с ключевыми идеями, мотивацией и проблемами обучения с подкреплением, а также помещает его в контекст более широкого пространства парадигм машинного обучения.


% ===========================================================
\section{Что такое обучение с подкреплением?}
\label{sec:what-is-rl}

\begin{definition}[Обучение с подкреплением]
\label{def:rl}
\index{reinforcement learning}
\textbf{Обучение с подкреплением} — это вычислительный подход к обучению на основе взаимодействия, при котором \emph{агент} совершает действия в \emph{среде}, получает \emph{вознаграждения} и стремится максимизировать суммарное вознаграждение на горизонте взаимодействия.
\end{definition}

Определяющая черта RL состоит в том, что агент никогда не получает явных указаний о том, какое действие является «правильным».  Вместо этого он получает лишь \emph{оценочную обратную связь}\index{evaluative feedback}: скалярный сигнал, указывающий, насколько хорошим (или плохим) оказался результат его поведения.  Задача агента, таким образом, заключается не в имитации заданной разметки, а в \emph{открытии стратегии}\,---\,называемой \emph{политикой}\index{policy}\,---\,которая максимизирует долгосрочную эффективность.

Чтобы формализовать качество поведения агента, введём понятие \emph{дохода}\index{return} (также называемого \emph{дисконтированным суммарным вознаграждением}), агрегирующего будущие вознаграждения начиная с шага $t$:
\begin{equation}
\label{eq:return-ch1}
G_t = \sum_{k=0}^{\infty}\gamma^k R_{t+k+1},
\end{equation}
где $\gamma\in[0,1)$ — \emph{коэффициент дисконтирования}\index{discount factor}, а $R_{t+k+1}$ — вознаграждение, полученное через $k+1$ шаг после момента $t$.  Следуя стандартному соглашению, на шаге~$t$ агент выбирает действие~$A_t$, после чего среда возвращает вознаграждение~$R_{t+1}$ и следующее состояние~$S_{t+1}$.

Если политика агента обозначена $\pi$, стандартная цель RL такова:
\begin{equation}
\label{eq:rl-objective}
J(\pi)=\E_{\pi}[G_0].
\end{equation}
Иными словами, мы ищем политику $\pi$, максимизирующую ожидаемый доход~\eqref{eq:rl-objective}.

\begin{remark}
\label{rem:episodic-continuing}
В \emph{эпизодических задачах}\index{episodic task} сумма в~\eqref{eq:return-ch1} обрывается в терминальном состоянии.  В \emph{непрерывных задачах}\index{continuing task} горизонт бесконечен, и коэффициент дисконтирования $\gamma<1$ обеспечивает конечность суммы.
\end{remark}


% ===========================================================
\section{Ключевые проблемы обучения с подкреплением}
\label{sec:rl-challenges}

Даже в простейших постановках RL возникает ряд фундаментальных проблем:

\begin{itemize}[leftmargin=*]
\item \textbf{Отложенное вознаграждение.}\index{delayed reward}
Последствия действия могут не проявиться вплоть до многих шагов спустя.  Ход в шахматах, кажущийся безобидным в дебюте, может определить исход партии пятьюдесятью ходами позже.

\item \textbf{Исследование против эксплуатации.}\index{exploration versus exploitation}
Агент должен находить баланс между \emph{исследованием}\,---\,пробой незнакомых действий с целью открыть потенциально лучшие стратегии\,---\,и \emph{эксплуатацией}\,---\,повторением действий, которые уже кажутся выгодными.  Это противоречие является центральной темой главы~\ref{ch:bandits}.

\item \textbf{Сдвиг распределения, вызванный политикой.}\index{distribution shift}
В отличие от обучения с учителем, где обучающее распределение обычно фиксировано, в RL собственная политика агента определяет, какие состояния и переходы он будет наблюдать далее.  Изменение политики меняет данные.

\item \textbf{Коррелированные данные.}
Последовательные наблюдения вдоль траектории временно коррелированы и не являются независимо взятыми, что усложняет применение стандартных гарантий статистического обучения.
\end{itemize}

\begin{remark}
Именно зависимость данных от текущей политики делает RL существенно более чувствительным к ошибкам на ранних этапах обучения.  Плохая начальная стратегия может не допустить агента в области пространства состояний-действий, где могло бы быть обнаружено более качественное поведение.
\end{remark}


% ===========================================================
\section{Краткая история обучения с подкреплением}
\label{sec:history}
\index{history of RL}

Интеллектуальные корни RL лежат на пересечении нескольких областей:

\begin{description}[leftmargin=!,labelwidth=3.2cm]
\item[Психология (1900-е--1950-е)]
\emph{Закон эффекта} Эдварда Торндайка (1911)\index{Law of Effect} гласил, что действия, за которыми следуют удовлетворяющие последствия, имеют тенденцию повторяться.  Эта идея обучения по сигналу вознаграждения является концептуальным предшественником RL.

\item[Оптимальное управление (1950-е)]
Ричард Беллман сформулировал теорию \emph{динамического программирования}\index{dynamic programming} и ввёл \emph{уравнение Беллмана}\index{Bellman equation}, которое по сей день остаётся основой алгоритмов RL.

\item[Стохастическая аппроксимация]
Роббинс и Монро (1951) разработали итерационные методы нахождения корней стохастических уравнений, заложив математический фундамент для инкрементных правил обучения.

\item[Временное разностное обучение (1980-е--1990-е)]
Алгоритм TD($\lambda$) Саттона (1988)\index{TD($\lambda$)} объединил идеи метода Монте-Карло и динамического программирования.  Уоткинс ввёл Q-learning (1989)\index{Q-learning}, а TD-Gammon Тесауро (1992)\index{TD-Gammon} продемонстрировал, что RL способен достигать уровня эксперта в нарды.

\item[Глубокий RL (2013 — настоящее время)]
Сочетание глубоких нейронных сетей с RL привело к серии прорывов: DQN играет в игры Atari на уровне человека (Mnih et al., 2013, 2015)\index{DQN}, AlphaGo обыгрывает чемпиона мира по го (Silver et al., 2016, 2017)\index{AlphaGo}, а алгоритмы управления с непрерывными действиями, такие как DDPG, SAC и PPO, обеспечивают манипуляции в робототехнике.

\item[RL для языковых моделей (2020 — настоящее время)]
Применение методов RL\,---\,в особенности PPO-based RLHF\index{RLHF}\,---\,для выравнивания больших языковых моделей (Ouyang et al., 2022; OpenAI, 2023) сделало RL центральным инструментом современного ИИ.  Методы DPO\index{DPO} (Rafailov et al., 2023) и GRPO\index{GRPO} (Shao et al., 2024) продолжают расширять границу RL--LLM.
\end{description}


% ===========================================================
\section{Обучение с подкреплением и другие парадигмы машинного обучения}
\label{sec:rl-vs-ml}
\index{supervised learning}\index{unsupervised learning}\index{imitation learning}

Чтобы оценить отличительные особенности RL, полезно сравнить его с другими ключевыми парадигмами машинного обучения.

\begin{table}[ht]
\centering
\renewcommand{\arraystretch}{1.3}
\begin{tabular}{@{}p{2.6cm}p{2.9cm}p{3.2cm}p{3.6cm}@{}}
\toprule
\textbf{Парадигма} & \textbf{Данные} & \textbf{Цель} & \textbf{Ключевая проблема} \\
\midrule
Обучение с учителем &
пары вход--метка $(x,y)$ &
минимизация ошибки предсказания &
обобщение, шум меток, сдвиг распределения \\
\addlinespace
Обучение без учителя &
неразмеченные данные $x$ &
обнаружение структуры, плотности или представлений &
выбор метрики качества; интерпретация \\
\addlinespace
Обучение с имитацией &
демонстрации эксперта &
воспроизведение политики эксперта &
сдвиг ковариат; накапливающиеся ошибки вдоль траекторий \\
\addlinespace
Обучение с подкреплением &
траектории взаимодействия со средой &
максимизация ожидаемого дохода &
исследование, отложенное вознаграждение, нестабильность обучения \\
\bottomrule
\end{tabular}
\caption{Сравнение обучения с подкреплением с другими основными парадигмами машинного обучения.}
\label{tab:rl-vs-ml}
\end{table}

Наиболее существенное отличие состоит в том, что в RL обучающее распределение \emph{не фиксировано}: собственная политика агента определяет, какие данные он будет собирать далее.  При изменении политики меняется и распределение наблюдений.

С практической точки зрения RL зачастую более требователен, чем обучение с учителем, по нескольким причинам:
\begin{itemize}[leftmargin=*]
\item Сбор данных дорогостоящ: он требует симуляции, взаимодействия с физической системой или большого числа эпизодов.
\item Сигнал вознаграждения нередко зашумлён, разрежён или плохо сконструирован.
\item Использование нейросетевых аппроксиматоров функций привносит нестабильности, требующие специализированных методов стабилизации (см. «смертельную триаду» в главе~\ref{ch:value-approx}).
\end{itemize}


% ===========================================================
\section{Применения обучения с подкреплением}
\label{sec:applications}
\index{applications}

Обучение с подкреплением было успешно применено в самых разнообразных областях:

\begin{description}[leftmargin=!,labelwidth=3.5cm]
\item[Игры]
Настольные игры (го, шахматы, сёги посредством AlphaZero), видеоигры (Atari, StarCraft~II, Dota~2), карточные игры (покер посредством Libratus, Pluribus).

\item[Робототехника]
Локомоция, ловкие манипуляции и перенос из симуляции в реальность с использованием PPO, SAC и модельных методов.

\item[Рекомендательные системы]
Оптимизация долгосрочной вовлечённости пользователей вместо мгновенного показателя кликабельности.

\item[Автономное вождение]
Сквозное обучение политик вождения на основе симуляции и реального взаимодействия.

\item[Исследование операций]
Управление запасами, динамическое ценообразование, распределение ресурсов и комбинаторная оптимизация.

\item[Здравоохранение]
Персонализированные стратегии лечения, адаптивные клинические испытания и оптимизация дозирования препаратов.

\item[Выравнивание языковых моделей]
RLHF для следования инструкциям, безопасности и полезности; DPO, GRPO и связанные методы для выравнивания LLM с человеческими предпочтениями.  Данному применению посвящена Часть~III этой книги.
\end{description}


% ===========================================================
\section{Интерфейс агент--среда}
\label{sec:agent-env}
\index{agent--environment interface}

На самом высоком уровне каждая задача RL включает два участника:

\begin{enumerate}
\item \textbf{Агент}\index{agent}: обучающийся и принимающий решения.
\item \textbf{Среда}\index{environment}: всё, с чем взаимодействует агент.
\end{enumerate}

На каждом дискретном шаге $t=0,1,2,\dots$ повторяется следующий цикл:
\begin{enumerate}
\item Агент наблюдает текущее состояние $S_t$ (или наблюдение $O_t$ в частично наблюдаемых постановках).
\item Агент выбирает действие $A_t$ в соответствии со своей политикой $\pi$.
\item Среда переходит в новое состояние $S_{t+1}$ и выдаёт скалярное вознаграждение $R_{t+1}$.
\end{enumerate}

Это взаимодействие порождает \emph{траекторию}\index{trajectory}:
\begin{equation}
\label{eq:trajectory}
S_0, A_0, R_1, S_1, A_1, R_2, S_2, A_2, R_3, \dots
\end{equation}

Формальная математическая рамка этого взаимодействия\,---\,\emph{марковский процесс принятия решений} (MDP)\,---\,является предметом главы~\ref{ch:mdp}.

\begin{figure}[ht]
\centering
\begin{tikzpicture}[
  block/.style={draw, rounded corners, minimum width=2.2cm, minimum height=1cm, align=center},
  arr/.style={-{Stealth[length=3mm]}, thick},
  every node/.style={font=\small}
]
\node[block] (agent) {Агент};
\node[block, right=4cm of agent] (env) {Среда};
\draw[arr] ([yshift= 6pt]agent.east) -- node[above]{$A_t$} ([yshift= 6pt]env.west);
\draw[arr] ([yshift=-3pt]env.west) -- node[below]{$R_{t+1},\; S_{t+1}$} ([yshift=-3pt]agent.east);
\end{tikzpicture}
\caption{Цикл взаимодействия агента и среды.  На каждом шаге агент посылает действие $A_t$ и получает вознаграждение $R_{t+1}$ и следующее состояние $S_{t+1}$.}
\label{fig:agent-env}
\end{figure}


% ===========================================================
\section{Обзор книги}
\label{sec:book-overview}

Остальная часть книги организована в четыре части.

\paragraph{Часть~I: Основы (главы~\ref{ch:bandits}--\ref{ch:td}).}
Мы начинаем с задачи многорукого бандита (глава~\ref{ch:bandits}), которая вводит компромисс между исследованием и эксплуатацией в чистейшем виде.  Затем формализуем полную постановку RL через марковские процессы принятия решений (глава~\ref{ch:mdp}), разрабатываем точные методы решения на основе динамического программирования (глава~\ref{ch:dp}) и вводим два семейства алгоритмов, основанных на выборке: методы Монте-Карло (глава~\ref{ch:mc}) и временное разностное обучение (глава~\ref{ch:td}).

\paragraph{Часть~II: Аппроксимация функций и глубокий RL (главы~\ref{ch:value-approx}--\ref{ch:actor-critic}).}
Мы рассматриваем проблему масштабирования RL на большие пространства состояний и действий посредством аппроксимации функций (глава~\ref{ch:value-approx}), методов градиента политики (глава~\ref{ch:pg}) и архитектур актор-критик, венчающихся PPO (глава~\ref{ch:actor-critic}).

\paragraph{Часть~III: RL для языковых моделей (главы~\ref{ch:reward-models}--\ref{ch:practical-rlhf}).}
Мы применяем инструментарий RL к обучению и выравниванию больших языковых моделей: начинаем с формирования вознаграждения и моделей вознаграждения (глава~\ref{ch:reward-models}), затем переходим к пайплайну дообучения LLM (глава~\ref{ch:lm-post-training}), методам DPO и прямого выравнивания (глава~\ref{ch:dpo}), GRPO и онлайн-RL с верифицируемыми вознаграждениями (глава~\ref{ch:grpo}), а также практике RLHF (глава~\ref{ch:practical-rlhf}).

\paragraph{Часть~IV: Направления исследований (главы~\ref{ch:reasoning}--\ref{ch:open}).}
Исследуем рассуждение и дообучение с подкреплением (глава~\ref{ch:reasoning}), мультиагентные системы (глава~\ref{ch:agents}) и завершаем обзором открытых проблем и перспективных направлений исследований (глава~\ref{ch:open}).


% ===========================================================
\section*{Упражнения}
\addcontentsline{toc}{section}{Упражнения}

\begin{exercise}
\label{ex:rl-examples}
Приведите три реальные задачи, которые естественно формулируются как задачи обучения с подкреплением.  Для каждой опишите агента, среду, пространство действий и сигнал вознаграждения.
\end{exercise}

\begin{exercise}
\label{ex:sl-vs-rl}
Рассмотрите задачу обучения робота ходьбе.  Объясните, почему подход с обучением с учителем (например, обучение по заранее записанным траекториям ходящего эксперта) может оказаться неудачным и почему подход RL может быть более подходящим.  Какую роль играет компромисс между исследованием и эксплуатацией?
\end{exercise}

\begin{exercise}
\label{ex:discount}
Покажите, что при $\gamma\in[0,1)$ и ограниченных вознаграждениях $|R_t|\le R_{\max}$ для всех $t$, доход~\eqref{eq:return-ch1} удовлетворяет
\[
|G_t| \le \frac{R_{\max}}{1-\gamma}.
\]
\end{exercise}

\begin{exercise}
\label{ex:paradigm-compare}
Для каждой из следующих задач определите, какая парадигма машинного обучения (с учителем, без учителя, с имитацией или с подкреплением) является наиболее естественной, и обоснуйте свой ответ:
\begin{enumerate}[label=(\alph*)]
\item Классификация электронных писем на спам и не-спам.
\item Обучение чат-бота быть полезным и безвредным с использованием данных о человеческих предпочтениях.
\item Обнаружение клиентских сегментов по истории покупок.
\item Обучение беспилотного автомобиля навигации по записям водителей-экспертов.
\end{enumerate}
\end{exercise}
